% ============================================================================
% Section 7 -- related work. Claim C4 and the novelty argument.
%
% STATUS: first full draft, 2026-09-05. Survives every branch of the GH200 gate.
% Source: seven papers read end-to-end 2026-08-10; positioning table in
% paper/PAPER_OUTLINE.md section 4 (that file is RETIRED as of
% 2026-09-06; the original is in git at 2bfc6ca5e).
%
% The endpoint argument is the load-bearing paragraph. Reviewers will go
% straight to "why does DirectKV not pre-empt this" -- answer it early and in
% their own words, with their own quotes.
% ============================================================================

\section{Related work}
\label{sec:related}

\paragraph{Capacity, not compute, is what binds.} A recent trace-driven study
of production serving finds that ``VRAM capacity, not token budget'' limits
batch size~\cite{bottlenecks}, and names scheduler improvements as future work.
That is the premise this paper builds on, and it is measured rather than
assumed.

\paragraph{Mechanisms for host-resident KV.} DirectKV~\cite{directkv}
contributes zero-copy GPU attention over CPU-resident KV on coherent
GH200/GB200 interconnects, through CPU-memory-aware tiling, warp pipelining and
kernel fusion, eliminating the GPU staging buffer entirely. It is the state of
the art in \emph{mechanism}, and it is our substrate rather than our competitor
--- its own \S8.1 says so: ``DirectKV does not change scheduling, batching, or
KV-management policies; instead, it provides an efficient execution path for
CPU-resident KV blocks\ldots while leaving the scheduler, model execution
graph, and frontend API unchanged.''

It is also worth recording that DirectKV's own analysis agrees with our
measured PCIe result: they conclude that ``on PCIe, DirectKV should be viewed
primarily as a capacity-extension mechanism\ldots its performance advantage
remains bounded by CPU--GPU interconnect bandwidth,'' and report NVLink-C2C
reducing attention latency by up to $4.2\times$ relative to PCIe. Our
\S\ref{sec:eval} finds the same boundary from the scheduling side. Their
evaluation covers Llama-3.1-8B and OPT-13B/30B on a GH200-96, so it does not
overlap the models we evaluate, and it does not include a 70B-class
configuration --- which our capacity model identifies as the most favourable
point of all, and which we therefore report as a confirmatory measurement
(\S\ref{sec:eval}).

\paragraph{The closest prior mechanism: partitioning by layer.}
Pie~\cite{pie} is the nearest neighbour of our transfer path and deserves a
precise comparison. It targets the same hardware, stages KV through GPU buffers
as we do, and prefetches per layer through a FIFO queue so that the transfer
overlaps compute --- our rolling region and its swap buffers are cousins, and it
reports up to $1.9\times$ throughput over vLLM.

\textbf{The difference is the axis of partition.} Pie relocates $m$ of a model's
$n$ \emph{layers} of KV to host memory, ``using per-layer KV cache as the basic
management unit'', and that split applies to \emph{every request identically};
its single control is a global expansion factor, tuned adaptively so that
swap time stays inside compute time. Two consequences follow directly. First,
\textbf{no request's KV is ever fully absent from the device} --- every request
retains $(n-m)/n$ of its cache in HBM, always, so there is no configuration in
which one request is entirely host-resident while its neighbours are entirely
device-resident. Second, \textbf{interconnect bandwidth is spent uniformly}: the
per-step volume is a fixed fraction of \emph{every} request's KV, whereas ours
is the KV of the requests we select, which is what makes \emph{which requests}
a decision at all.

\paragraph{A consequence of the layer unit: the return leg.} Choosing the layer
as the unit of management has a second effect, on traffic rather than on
policy. In Pie a layer's KV \emph{lives} on one side or the other, and the
GPU-side buffer it occupies is needed by the next layer arriving, so a layer
that has just been computed --- and has therefore just had this token's $K,V$
appended to it --- must be written back in full. Their swap controller says so
directly: it initiates ``a swap-out for the `coldest layer', defined as the last
layer that finished computation''. Every offloaded byte therefore crosses the
interconnect \textbf{twice per generated token}, and the return leg is the
slower direction (371 vs 419\,GB/s by their own measurement).

Our rolling region is not the home of the data but a cache of it: the host
mirror stays authoritative and each step's new $K,V$ is written to it directly,
so a slot can be overwritten without persisting anything. DirectKV's kernel has
the same property, appending only the newly projected $K,V$ to CPU memory. The
resulting difference in per-token traffic is not marginal:

\begin{table}[t]
\centering
\small
\begin{tabular}{lll}
\toprule
per generated token & to device & to host \\
\midrule
Pie~\cite{pie} & offloaded KV & \textbf{the same KV again} \\
DirectKV~\cite{directkv} & all KV & one token's $K,V$ \\
\textbf{this paper} & selected requests' KV & one token's $K,V$ \\
\bottomrule
\end{tabular}
\caption{Per byte of offloaded KV, a design whose offloaded unit is \emph{owned}
by one side must return it; one whose device-side copy is a \emph{cache} need
not. Pie moves two bytes where the other two move one, and its return leg uses
the slower direction of the link.}
\label{tab:traffic}
\end{table}

This is not a criticism of Pie's design so much as a property of its unit:
a whole layer's cache cannot be left dirty on the device when the buffer is
reclaimed for the next layer. It does mean that comparisons of interconnect
pressure between these systems must count both directions.

The two axes are orthogonal rather than competing, and we say so deliberately: a
layer-split of the kind Pie performs could operate \emph{within} our device
tier. We also adopt one of its measurements --- Pie reports GH200 NVLink-C2C at
419\,GB/s CPU$\to$GPU and 371\,GB/s GPU$\to$CPU, an independent confirmation
that the widely quoted 900\,GB/s is a bidirectional figure and that a
host-to-device KV stream should be budgeted well below it.

FlexGen~\cite{flexgen} is the earlier generation of the same family, placing KV
across GPU, CPU and disk by a linear-programming search with compression; it is
offline-profiled rather than adaptive. Both pay the staging cost that a
zero-copy kernel removes, and neither decides a tier per request.

DAK~\cite{dak} pushes host-resident tensors to shared memory via TMA and chooses
per-\emph{operation} offload ratios by greedy search --- but takes a
\emph{global} offload ratio as an input, which is precisely the quantity a
policy has to produce.

\paragraph{The endpoint argument.} DirectKV \S6: ``During prefill, the manager
stores the initial $K,V$ entries from each layer in CPU memory. During decode,
the GPU reuses these cached entries directly, while only appending new $K,V$
for the latest token'' --- and its Alg.~3 appends that new $K,V$ to CPU memory
as well. \textbf{It places 100\% of KV on the host, for every request, always.}
That is why it needs no scheduler: placement is static and total. Stock vLLM and SGLang are the opposite endpoint, at 0\%.
Nobody runs the mixture, and the optimum is interior
(Table~\ref{tab:endpoints}). Note in particular that the all-host endpoint
leaves the device memory not holding weights entirely idle, which is why it
loses to an HBM-resident baseline whenever that baseline fits --- a result
DirectKV reports on Llama-3.1-8B and does not explain, and which our gain law
predicts (\S\ref{sec:model:directkv}).

A natural objection to the composability reading of DirectKV \S8.1 is that the
vLLM and SGLang mechanisms it names --- continuous batching, prefix reuse,
eviction, hierarchical KV caching --- already constitute KV policy. They do
not, for our purposes: each decides \emph{reuse} and \emph{preemption}, that
is, the residency of blocks belonging to a request that is \emph{not currently
decoding}. None of them decides the tier of an \emph{actively decoding}
request. That state did not exist before the zero-copy mechanism created it.

\paragraph{Policy, at other granularities.} OmniServe~\cite{omniserve} is the
closest prior work on policy: it forms mixed GPU/host batches and admits
host-tier work per layer under an SLO. The regime differs on three axes that
change the problem. It targets PCIe-attached A100s with CPU-only hosts over
RoCE; attention for host-tier requests is computed \emph{on the CPU}, so the
trade is compute-for-capacity rather than bandwidth-for-capacity; and tiering
is driven by an \emph{externally supplied} latency-sensitive/best-effort
service class. Ours assumes a coherent interconnect, computes all attention on
the GPU, and infers tier from observable request properties, so it applies in a
single-class deployment where no such label exists.

Dynamic KV Placement~\cite{calkv} gives a formal placement model and a
simulated-annealing oracle bound on aggregate bandwidth. It is simulation only,
on synthetically shrunk HBM, at \emph{token} granularity by importance, and it
requires perfect foresight; its gains are contingent on 60--80\% attention
sparsity, and it closes by explicitly calling for runtime policies. Its own
results show migration overhead dominating below roughly 60\% sparsity. We
make the testable counter-claim that \emph{request}-granularity migration
avoids that thrash, because whole requests move once at controlled moments
rather than pages chasing shifting token importance.

\paragraph{Reducing the KV instead of relocating it.} ShadowKV~\cite{shadowkv}
keeps a low-rank $K$ on the GPU, offloads $V$, and reconstructs sparsely,
reporting a $6\times$ footprint reduction and $3.04\times$ throughput; KV
quantization, Quest and H2O likewise shrink what must be stored. These change
the model's output, and they compose with a memory system rather than competing
with it: a smaller KV scales both the device-resident cohort and the affordable
host cohort, leaving the ratio between them --- and hence the value of a second
tier --- unchanged. \S\ref{sec:eval} measures that composition on KV
quantization rather than assuming it.

\paragraph{Moving compute instead of data.} NEO~\cite{neo} offloads attention
\emph{compute} to the CPU to grow the batch --- the same problem statement as
ours with the opposite mechanism. Its own scaling is the argument against it:
gains decay from $7.5\times$ on a T4 to 14\% on an H100, because a faster GPU
makes the CPU's arithmetic relatively worse. Streaming KV to the GPU has the
opposite dependence: it improves as the interconnect improves, which is the
direction hardware is moving.

\paragraph{Summary of the gap.} Table~\ref{tab:axes} places the systems that
share our hardware assumption on the two questions that matter: what unit of KV
is placed, and who decides.

\begin{table}[t]
\centering
\small
\begin{tabular}{llll}
\toprule
 & placement unit & decided by & device KV \\
\midrule
vLLM, SGLang & --- & --- & all of it \\
Pie~\cite{pie} & layer, uniform & global factor & $(n{-}m)/n$ each \\
DirectKV~\cite{directkv} & none (all host) & nothing & none \\
\textbf{this paper} & \textbf{request} & \textbf{placement,} & resident \\
 & & \textbf{admission,} & cohort \\
 & & \textbf{migration} & only \\
\bottomrule
\end{tabular}
\caption{The mechanism for reading host-resident KV efficiently is solved; the
open question is which requests should live where. Systems either place nothing
(both endpoints) or place a fixed fraction of everything (Pie), and none decides
per request.}
\label{tab:axes}
\end{table}

What is unsolved is which requests should live in which tier, how many host-tier
requests a step can absorb without breaking the resident cohort's latency, and
when a request should move. Three of the papers above name that layer as out of
scope or future work.
